\textbf{Phần 1. Trắc nghiệm lựa chọn}
\setcounterpageref{ex}{0}

\begin{ex} %[5D4V23]
Tính giá trị của biểu thức: $4{,}85 \times 100 - 325 \times 0{,}1$.
\choice
{$451{,}5$}
{\True $452{,}5$}
{$453{,}5$}
{$462{,}5$}
\loigiai{
Ta có:
$$4{,}85 \times 100 - 325 \times 0{,}1 = 485 - 32{,}5 = 452{,}5.$$
\textbf{Đáp án của câu hỏi là B}
}
\end{ex}

\begin{ex} %[5D4V23]
Khi dời dấu phẩy của một số thập phân sang bên phải một chữ số, ta được một số mới lớn hơn số ban đầu $16{,}38$ đơn vị. Số thập phân ban đầu là
\choice
{$18{,}2$}
{\True $1{,}82$}
{$0{,}182$}
{$1{,}638$}
\loigiai{
Khi dời dấu phẩy sang phải một chữ số thì số đó gấp lên $10$ lần, nghĩa là tăng thêm $9$ lần số ban đầu.\\
Số thập phân ban đầu là:
$16{,}38 : 9 = 1{,}82$.\\
\textbf{Đáp án của câu hỏi là B}
}
\end{ex}

\begin{ex} %[5D4V23]
Khi dời dấu phẩy của một số thập phân sang bên trái một chữ số, ta được số mới bé hơn số ban đầu $29{,}25$ đơn vị. Tổng của số mới và số ban đầu là
\choice
{$32{,}5$}
{\True $35{,}75$}
{$34{,}5$}
{$36{,}25$}
\loigiai{
Khi dời dấu phẩy sang trái một chữ số thì số ban đầu gấp $10$ lần số mới.\\
Số mới là:
$29{,}25 : (10 - 1) = 3{,}25$.\\
Số ban đầu là:
$3{,}25 \times 10 = 32{,}5$.\\
Tổng của hai số là:
$32{,}5 + 3{,}25 = 35{,}75$.\\
\textbf{Đáp án của câu hỏi là B}
}
\end{ex}

\begin{ex} %[5D4V23]
Tính thuận tiện: $7{,}8 \times 10 + 780 \times 0{,}1 + 0{,}78 \times 100$.
\choice
{$78$}
{$156$}
{\True $234$}
{$312$}
\loigiai{
Ta có:
$$7{,}8 \times 10 = 78;\quad 780 \times 0{,}1 = 78;\quad 0{,}78 \times 100 = 78.$$
Tổng bằng:
$$78 + 78 + 78 = 234.$$
\textbf{Đáp án của câu hỏi là C}
}
\end{ex}

\begin{ex} %[5D4V23]
Tìm $x$, biết: $x : 0{,}01 = 250 \times 0{,}1 + 75$.
\choice
{$0{,}1$}
{\True $1$}
{$10$}
{$100$}
\loigiai{
Ta có:
\begin{align*}
x : 0{,}01 &= 25 + 75\\
x : 0{,}01 &= 100\\
x &= 100 \times 0{,}01 = 1
\end{align*}
\textbf{Đáp án của câu hỏi là B}
}
\end{ex}

\begin{ex} %[5D4V23]
Một tấm bìa hình chữ nhật có chiều rộng $0{,}12\text{ m}$. Chiều dài gấp $10$ lần chiều rộng. Diện tích của tấm bìa đó tính theo đề-xi-mét vuông là
\choice
{$1{,}44\text{ dm}^2$}
{\True $14{,}4\text{ dm}^2$}
{$144\text{ dm}^2$}
{$0{,}144\text{ dm}^2$}
\loigiai{
Đổi: $0{,}12\text{ m} = 1{,}2\text{ dm}$.\\
Chiều dài của tấm bìa là:
$1{,}2 \times 10 = 12\text{ (dm)}$.\\
Diện tích tấm bìa là:
$12 \times 1{,}2 = 14{,}4\text{ (dm}^2\text{)}$.\\
\textbf{Đáp án của câu hỏi là B}
}
\end{ex}

\begin{ex} %[5D4V23]
Một số khi nhân với $0{,}1$ rồi cộng thêm $15{,}4$ thì được kết quả là $18{,}2$. Số đó là
\choice
{$2{,}8$}
{\True $28$}
{$280$}
{$0{,}28$}
\loigiai{
Tích của số đó với $0{,}1$ là:
$18{,}2 - 15{,}4 = 2{,}8$.\\
Số cần tìm là:
$2{,}8 : 0{,}1 = 28$.\\
\textbf{Đáp án của câu hỏi là B}
}
\end{ex}

\begin{ex} %[5D4V23]
Tính nhanh giá trị của biểu thức $2024 \times 0{,}1 + 2024 \times 0{,}9$:
\choice
{$202{,}4$}
{$20{,}24$}
{\True $2\,024$}
{$20\,240$}
\loigiai{
Đặt thừa số chung:
$$2024 \times (0{,}1 + 0{,}9) = 2024 \times 1 = 2\,024.$$
\textbf{Đáp án của câu hỏi là C}
}
\end{ex}

\begin{ex} %[5D4V23]
Một bao gạo nặng $50\text{ kg}$. Lần đầu người ta lấy ra $0{,}1$ khối lượng gạo trong bao, lần sau lấy tiếp $0{,}2$ khối lượng gạo còn lại. Hỏi sau hai lần lấy, trong bao còn lại bao nhiêu ki-lô-gam gạo?
\choice
{$35\text{ kg}$}
{\True $36\text{ kg}$}
{$40\text{ kg}$}
{$38\text{ kg}$}
\loigiai{
Lần đầu lấy ra số ki-lô-gam gạo là:
$50 \times 0{,}1 = 5\text{ (kg)}$.\\
Số gạo còn lại sau lần một là:
$50 - 5 = 45\text{ (kg)}$.\\
Lần hai lấy ra số ki-lô-gam gạo là:
$45 \times 0{,}2 = 9\text{ (kg)}$.\\
Sau hai lần lấy, trong bao còn lại:
$45 - 9 = 36\text{ (kg)}$.\\
\textbf{Đáp án của câu hỏi là B}
}
\end{ex}

\begin{ex} %[5D4V23]
Phép tính nào sau đây cho kết quả khác với các phép tính còn lại?
\choice
{$45 \times 0{,}1$}
{$450 : 100$}
{$0{,}45 \times 10$}
{\True $4{,}5 : 0{,}1$}
\loigiai{
Ta có:\\
$45 \times 0{,}1 = 4{,}5$;\\
$450 : 100 = 4{,}5$;\\
$0{,}45 \times 10 = 4{,}5$;\\
Nhưng $4{,}5 : 0{,}1 = 45 
eq 4{,}5$.\\
\textbf{Đáp án của câu hỏi là D}
}
\end{ex}

\begin{ex} %[5D4V23]
Dời dấu phẩy của số thập phân $A$ sang bên phải hai chữ số ta được số thập phân $B$. Biết rằng $A + B = 254{,}52$. Số thập phân $A$ là
\choice
{$25{,}2$}
{\True $2{,}52$}
{$0{,}252$}
{$2{,}54$}
\loigiai{
Khi dời dấu phẩy sang phải hai chữ số thì số $B$ gấp $100$ lần số $A$.\\
Do đó:
$A + B = A + 100 \times A = 101 \times A = 254{,}52$.\\
Số $A$ là:
$254{,}52 : 101 = 2{,}52$.\\
\textbf{Đáp án của câu hỏi là B}
}
\end{ex}

\begin{ex} %[5D4V23]
Một trang trại có $1\,000$ cây ăn quả. Trong đó số cây cam chiếm $0{,}45$ tổng số cây, số cây xoài chiếm $0{,}35$ tổng số cây, còn lại là cây bưởi. Số cây bưởi trong trang trại là
\choice
{$150\text{ cây}$}
{\True $200\text{ cây}$}
{$250\text{ cây}$}
{$300\text{ cây}$}
\loigiai{
Số cây bưởi chiếm số phần tổng số cây là:
$1 - (0{,}45 + 0{,}35) = 1 - 0{,}8 = 0{,}2$.\\
Số cây bưởi là:
$1\,000 \times 0{,}2 = 200\text{ (cây)}$.\\
\textbf{Đáp án của câu hỏi là B}
}
\end{ex}

\textbf{Phần 2. Trắc nghiệm đúng sai}
\setcounterpageref{ex}{0}

\begin{ex} %[5D4V23]
Cho ba biểu thức: $x = 0{,}65 \times 100$; $y = 650 \times 0{,}1$; $z = 6{,}5 : 0{,}1$.
\choiceTF
{\True Giá trị của $x$ bằng $65$}
{\True Giá trị của $y$ bằng $65$}
{\True Ba biểu thức $x, y, z$ có cùng giá trị}
{Tổng $x + y + z = 185$}
\loigiai{
\begin{itemchoice}
\item Đúng. $x = 0{,}65 \times 100 = 65$.
\item Đúng. $y = 650 \times 0{,}1 = 65$.
\item Đúng. $z = 6{,}5 : 0{,}1 = 65$, do đó $x = y = z = 65$.
\item Sai. Tổng $x + y + z = 65 \times 3 = 195 
eq 185$.
\end{itemchoice}
}
\end{ex}

\begin{ex} %[5D4V23]
Một cuộn dây thép dài $240\text{ m}$. Người ta cắt lấy $0{,}1$ chiều dài sợi dây để làm móc phơi, sau đó lấy tiếp $0{,}25$ chiều dài còn lại để rào luống rau.
\choiceTF
{\True Chiều dài đoạn dây cắt lần thứ nhất là $24\text{ m}$}
{\True Chiều dài đoạn dây còn lại sau lần cắt thứ nhất là $216\text{ m}$}
{\True Chiều dài đoạn dây cắt lần thứ hai là $54\text{ m}$}
{Sau hai lần cắt, cuộn dây còn lại $150\text{ m}$}
\loigiai{
\begin{itemchoice}
\item Đúng. Chiều dài đoạn cắt lần 1: $240 \times 0{,}1 = 24\text{ (m)}$.
\item Đúng. Đoạn còn lại sau lần 1: $240 - 24 = 216\text{ (m)}$.
\item Đúng. Chiều dài đoạn cắt lần 2: $216 \times 0{,}25 = 54\text{ (m)}$.
\item Sai. Đoạn còn lại là: $216 - 54 = 162\text{ (m)} 
eq 150\text{ m}$.
\end{itemchoice}
}
\end{ex}

\textbf{Phần 3. Trắc nghiệm trả lời ngắn}
\setcounterpageref{ex}{0}

\begin{ex} %[5D4V23]
Tính giá trị của biểu thức: $12{,}5 \times 0{,}1 \times 100 \times 8$.
\shortans{1000}
\loigiai{
Ta nhóm các thừa số thích hợp:
$$(12{,}5 \times 8) \times (0{,}1 \times 100) = 100 \times 10 = 1\,000.$$
}
\end{ex}

\begin{ex} %[5D4V23]
Tìm số mới khi dời dấu phẩy của một số thập phân sang bên phải một chữ số, biết rằng tổng của số mới và số ban đầu bằng $132$.
\shortans{120}
\loigiai{
Số mới gấp 10 lần số ban đầu.\\
Tổng số phần bằng nhau là: $1 + 10 = 11$ (phần).\\
Số ban đầu là: $132 : 11 = 12$.\\
Số mới cần tìm là: $12 \times 10 = 120$.
}
\end{ex}

\begin{ex} %[5D4V23]
Tìm số tự nhiên $x$ chia hết cho $10$ thỏa mãn: $15 < x \times 0{,}1 < 17$.
\shortans{160}
\loigiai{
Chia cả biểu thức cho $0{,}1$ (tương đương nhân với $10$):
$$150 < x < 170.$$
Vì $x$ là số tự nhiên chia hết cho 10 nên $x = 160$.
}
\end{ex}

\begin{ex} %[5D4V23]
Một kho chứa $150$ tấn thóc. Ngày thứ nhất người ta xuất đi $0{,}1$ số thóc trong kho. Ngày thứ hai người ta xuất tiếp $0{,}2$ số thóc còn lại. Hỏi sau hai ngày trong kho còn lại bao nhiêu tấn thóc?
\shortans{108}
\loigiai{
Ngày thứ nhất xuất: $150 \times 0{,}1 = 15\text{ (tấn)}$.\\
Số thóc còn lại sau ngày thứ nhất: $150 - 15 = 135\text{ (tấn)}$.\\
Ngày thứ hai xuất: $135 \times 0{,}2 = 27\text{ (tấn)}$.\\
Sau hai ngày còn lại: $135 - 27 = 108\text{ (tấn)}$.
}
\end{ex}

\begin{ex} %[5D4V23]
Tính giá trị của biểu thức: $35 \times 0{,}1 + 45 \times 0{,}1 + 20 \times 0{,}1$.
\shortans{10}
\loigiai{
Đặt thừa số chung $0{,}1$:
$$(35 + 45 + 20) \times 0{,}1 = 100 \times 0{,}1 = 10.$$
}
\end{ex}

\begin{ex} %[5D4V23]
Dời dấu phẩy của một số thập phân sang bên trái một chữ số thì số đó giảm đi $57{,}6$ đơn vị. Tìm số thập phân ban đầu.
\shortans{64}
\loigiai{
Khi dời dấu phẩy sang trái một chữ số thì số ban đầu gấp 10 lần số mới, tức là giảm đi 9 lần số mới.\\
Số mới là:
$57{,}6 : 9 = 6{,}4$.\\
Số thập phân ban đầu là:
$6{,}4 \times 10 = 64$.
}
\end{ex}
